% =====================================================================
%  Chapter 14 opener -- "Chapter at a Glance", three reading paths + the
%  ladder staircase, THE running example (the same 10-node tree T as
%  Ch.12/Ch.13, edges 1-2,1-3,2-4,2-5,3-6,3-7,5-8,7-9,9-10) re-drawn with
%  a budgeted tree-knapsack DP table overlay (V=3), the "who owns what"
%  comparison (Ch.13 fixed-tuple cookbook vs THIS chapter's budgeted and
%  game DPs vs Ch.15 rerooting vs Ch.16 LCA vs Ch.17 DSU vs Ch.24 centroid),
%  roadmap, recognition guide, constraint->tool table, notation.
%  Exports fig:c27-runtree-knap and fig:c27-taxonomy.
%  All DP tables below are hand-computed on T; every later section re-uses
%  exactly these numbers.  Node value = node label, every item costs 1, so
%  dp[v][j] = sum of the j largest labels in subtree(v): a merge whose
%  RESULT is trivial to verify while its LOOP COUNT is what 14.1 counts.
% =====================================================================

\begin{chaptertheorybox}{Chapter at a Glance}
Chapter~13 taught the tree DP with a \emph{constant} tuple: one small bag of
numbers per node, folded from the children in $\bigO{\deg(v)}$, one DFS, done.
This chapter is what happens the moment the tuple acquires an
\term{index} -- a budget, a count, a length, a value cap -- so the state at a
node is an \emph{array} $\mathit{dp}[v][0\,.\,V]$ rather than a scalar.  The
recurrence no longer combines children independently: budget spent on one
child cannot be spent on the next, so children must be merged \emph{one at a
time} in a knapsack loop, and the naive cost analysis (the subtree sizes
multiply!) looks quadratic.  The chapter's spine is the proof that this
illusion is false:

\term{the amortised child-merge} (\secref{c27-knapsack}):
$\mathit{dp}[v][j] = \max_{\text{split}} \sum_c \mathit{dp}[c][j_c]$
with $\sum_c j_c \le j$; merging children in order, bounding the inner loops
by $\min(\text{size so far},V)\times\min(\text{child},V)$, the whole DP costs
$\bigO{nV}$ \emph{across all nodes} (\thmref{c27-amortize}) -- the bound Ch.~13
previewed twice and never proved.  CF~855C's cup-counting DP is the worked
example.
\term{states that grow} (\secref{c27-growing}): when the index is not a fixed
small $V$ but a value/time/length that ranges with the data, the state is a
\emph{profile}; the merge of 14.1 still keeps it linear, and CF~490F's
longest-increasing-subsequence-along-a-path is the case study.
\term{greedy where a DP would fit} (\secref{c27-greedy}): sometimes the
budgeted DP collapses to a single leaf-pruning sweep because a tree has no
cycles -- a cut-then-swap exchange argument does the work of the table.
CF~1029E's distance-$2$ cover is the proof-by-hand.
\term{games on trees} (\secref{c27-games}): two players, optimal play, and the
value of a position is a $\mathrm{mex}$ or a win/lose flag propagated up the
trie; CF~455B's four-state per-node DP is the paradigm, Sprague--Grundy the
general theory.

\textbf{You will learn to}
\begin{itemize}[nosep]
  \item write the budgeted child-merge template and \emph{bound its loops}
        by the running subtree size so the DP is $\bigO{nV}$, not $\bigO{n^2}$; \hfill\lv{2}--\lv{3}
  \item state and prove \thmref{c27-amortize} (the sum-of-subtree-size
        argument), Ch.~13's outstanding debt; \hfill\lv{4}
  \item drive a value/length-indexed profile DP cold on CF~490F; \hfill\lv{4}--\lv{5}
  \item recognise when a tree DP has a greedy shadow and prove it by
        cut-and-swap (CF~1029E); \hfill\lv{3}--\lv{4}
  \item compute Grundy numbers on a tree and read a win from a
        $\mathrm{mex}$/XOR, then solve the win/lose trie DP CF~455B; \hfill\lv{4}
  \item combine everything on the Summit CF~735E: a budgeted subtree knapsack
        ($f[v][j]$, $j\le$ distance-to-nearest-black) merged under the
        $\bigO{nV}$ cap. \hfill\lv{5}
\end{itemize}

\textbf{Prerequisites.} Chapter~13 (\chref{c10}) for the child-fold template
\thmref{c10-subtree-template} and the two histogram shapes this chapter budgets;
Chapter~12 (\chref{c09}) for the running tree $T$, the depth/size one-DFS
metrics and \thmref{c09-dpjoin}; Chapter~4 (\chref{c04}) for rooting by a
parent parameter. Small-to-large merging is \emph{deliberately not assumed} --
it is Chapter~17's (\chref{c28}); every merge here is a plain per-child knapsack.
\end{chaptertheorybox}

% ---------------- Reading paths + ladder position ----------------
\begin{chaptertheorybox}{How to Read This Chapter (three paths)}
Level badges $\lv1$--$\lv6$: $\lv2$ is the merge skeleton you write straight
from the recurrence; $\lv3$ is a solid budgeted DP coded from a section's
listing; $\lv4$ needs \emph{two} ideas composed (a growing index inside the
$\bigO{nV}$ amortisation, or a greedy justified by an exchange proof);
$\lv5$ is the Summit, where a distance budget and a subtree knapsack must
survive together; $\lv6$ is reading the amortised bound as a reusable tool.

\textbf{Amortisation-first path ($\lv2$--$\lv{4}$, ``I owe Ch.13 a proof'').}
Read \secref{c27-knapsack} end to end: the template, \thmref{c27-amortize},
then the CF~855C solved problem and its trace. You leave able to bound any
"merge children with a size/count budget" DP by $\bigO{nV}$ -- which is the
one technique Ch.~13 used but could not justify.

\textbf{Growing-states path ($\lv4$).} Start at \secref{c27-growing}: the
index is no longer a tiny constant but a value/time/length that ranges with
the data; carry the amortisation of \secref{c27-knapsack} as a black box and
drive CF~490F. This is the reading that turns Ch.~13's "tuple size blows up"
pitfall into a designed state.

\textbf{Games-and-greedy path ($\lv3$--$\lv{5}$, contest track).} Read
\secref{c27-greedy} and \secref{c27-games} (they need only \secref{c27-knapsack}'s
loop-bounding discipline), solve CF~1029E and CF~455B, then the mixed set,
then the Summit CF~735E. At that point the forward ropes are visible:
small-to-large is Chapter~17 (\chref{c28}), rerooting-as-theory is
Chapter~15 (\chref{c20}), pair distance \emph{queries} are Chapter~16 (\chref{c11}).

The backward ropes: Chapter~13 (\chref{c10}) for the fixed-tuple cookbook this
chapter budgets on, and Chapter~12 (\chref{c09}) for $T$ and the metrics.
\end{chaptertheorybox}

\begin{figure}[H]
\centering
\begin{tikzpicture}[diagram, x=9mm, y=7mm]
  \foreach \st/\lab in {1/lv1,2/lv2,3/lv3,4/lv4,5/lv5,6/lv6} {
    \fill[\lab!22, draw=\lab, line width=0.8pt] (\st-1,0) rectangle (\st,\st*0.52);
    \node[font=\scriptsize\bfseries, text=\lab] at (\st-0.5,\st*0.52-0.42) {L\st};}
  \node[annot, anchor=south west, font=\scriptsize]
    at (0.05,0.62) {14.1 budgeted child-merge skeleton};
  \node[annot, anchor=south west, font=\scriptsize]
    at (1.05,1.14) {14.1 \thmref{c27-amortize} loop cap $\min(\cdot,V)$};
  \node[annot, anchor=south west, font=\scriptsize]
    at (2.05,1.66) {14.1 CF 855C cup DP $\mathit{dp}[v][\text{cls}][\text{cnt}]$};
  \node[annot, anchor=south west, font=\scriptsize, text=lv5]
    at (3.05,2.18) {14.2 CF 490F profile DP; 14.3 CF 1029E exchange};
  \node[annot, anchor=south west, font=\scriptsize, text=lv5]
    at (4.05,2.70) {14.4 CF 455B win/lose trie; Grundy $\mathrm{mex}$/XOR};
  \node[annot, anchor=south west, font=\scriptsize, text=lv6]
    at (5.05,3.22) {Summit CF 735E budgeted subtree knapsack};
\end{tikzpicture}
\caption*{Chapter~14 on the six-step ladder: the merge skeleton and the
amortised loop cap sit at $\lv2$--$\lv3$; CF~855C's coloured cup-count DP and
the cut-and-swap greedy of CF~1029E are $\lv3$--$\lv4$; the value-indexed
profile of CF~490F and the four-state game DP of CF~455B climb to $\lv4$--$\lv5$;
the Summit CF~735E (a distance-budget subtree knapsack) is step~5 read from
step~6. Amortisation path $=$ steps 2--4; full contest path runs to 6.}
\end{figure}

% ---------------- The running example ----------------
\subsection*{The Running Example: One Ten-Node Tree, a Budget on the State}

\emph{One} rooted tree does all the work -- the same $T$ Chapters~12 and~13
used, now with a knapsack budget bolted on. Nodes $1,\dots,10$, root fixed at
$1$, edges
\[
  1\!-\!2,\; 1\!-\!3,\; 2\!-\!4,\; 2\!-\!5,\; 3\!-\!6,\; 3\!-\!7,\; 5\!-\!8,\; 7\!-\!9,\; 9\!-\!10 .
\]
Child sets in the rooted view: $1\!\to\!\{2,3\}$, $2\!\to\!\{4,5\}$,
$3\!\to\!\{6,7\}$, $5\!\to\!\{8\}$, $7\!\to\!\{9\}$, $9\!\to\!\{10\}$; leaves
$\{4,6,8,10\}$. Subtree sizes $\size(v)$ are $10,4,5,1,2,1,3,1,2,1$ for
$v=1..10$ (Ch.~13's table, reused verbatim), sum $=30$ -- the very quantity
\thmref{c27-amortize} bounds the merge by.

The overlay in \figref{c27-runtree-knap} is the \emph{budgeted knapsack}
$\mathit{dp}[v][j]$: pick a subset of $\text{subtree}(v)$ of total cost $\le j$,
each node costs~$1$ and is worth its label, maximise total value. Because every
item costs~$1$, the answer is just \emph{the sum of the $j$ largest labels in
$\text{subtree}(v)$} -- trivial to hand-check, which is the point: the RESULT
is easy to verify while the \emph{loop count} of the merge is what 14.1 is
really studying. Budget $V=3$.

\begin{figure}[H]
\centering
\begin{tikzpicture}[diagram, scale=0.90, transform shape,
    tnode/.style={circle, draw=primaryblue, line width=0.9pt, fill=primaryblue!8,
                  minimum size=8mm, font=\small\bfseries, text=primaryblue!60!black}]
  \node[tnode] (n1) at (0,0)     {$1$};
  \node[tnode] (n2) at (-3,-1.6) {$2$};
  \node[tnode] (n3) at ( 3,-1.6) {$3$};
  \node[tnode] (n4) at (-5.2,-3.2) {$4$};
  \node[tnode] (n5) at (-1.6,-3.2) {$5$};
  \node[tnode] (n6) at ( 1.6,-3.2) {$6$};
  \node[tnode] (n7) at ( 5.2,-3.2) {$7$};
  \node[tnode] (n8) at (-1.6,-4.8) {$8$};
  \node[tnode] (n9) at ( 5.2,-4.8) {$9$};
  \node[tnode] (n10) at ( 5.2,-6.4) {$10$};
  \draw[tedge] (n1)--(n2) (n1)--(n3) (n2)--(n4) (n2)--(n5)
               (n3)--(n6) (n3)--(n7) (n5)--(n8) (n7)--(n9) (n9)--(n10);
  % knapsack badge: dp[v][1] / dp[v][2] / dp[v][3] (top-j label sums)
  \foreach \nd/\tr in {n1/10-19-27, n2/8-13-17, n3/10-19-26, n4/4,
                      n5/8-13, n6/6, n7/10-19-26, n8/8, n9/10-19, n10/10} {
    \node[font=\tiny\ttfamily, text=warningred!75!black, anchor=north]
      at (\nd.south) {(\tr)};}
  \node[annot, font=\scriptsize, anchor=south west, text=successgreen!50!black]
    at (-0.4,0.55) {read dp[v][1..3] below each node};
\end{tikzpicture}
\caption{The running tree $T$ with the \emph{budgeted knapsack} overlay, $V=3$.
The triple below each node lists $\mathit{dp}[v][1],\mathit{dp}[v][2],\mathit{dp}[v][3]$ (top-$j$
label sums in $\text{subtree}(v)$; a leaf shows only $\mathit{dp}[\cdot][1]$, a
$2$-node subtree only up to $[\cdot][2]$). Sanity: root $\mathit{dp}[1][3]=27=10{+}9{+}8$,
the three largest labels; the merge of $\mathit{dp}[2][\cdot]$ and
$\mathit{dp}[3][\cdot]$ into $\mathit{dp}[1][\cdot]$ is the child-merge
\secref{c27-knapsack} bounds at $\bigO{nV}$.}
\label{fig:c27-runtree-knap}
\end{figure}

Four angles, four hand-computed overlays -- each lives in the section named on
its row and never changes:

\begin{center}
\small
\renewcommand{\arraystretch}{1.25}
\begin{tabular}{@{}p{0.26\linewidth} p{0.46\linewidth} l@{}}
\toprule
\textbf{Angle on $T$} & \textbf{What comes out} & \textbf{Traced in} \\
\midrule
knapsack merge cost & inner-loop iterations $\sum_v\sum_c \min(\text{so far},3)\!\cdot\!\min(\text{child},3)=28\le 3n$ (uncapped the same merge is $\sum_v\!\!\sum_c \text{so far}\!\cdot\!\text{child}$, far larger) & \secref{c27-knapsack} \\
profile / LIS-along-path & $\mathit{dp}$ indexed by a value cap, longest strictly-increasing concert sequence over a path & \secref{c27-growing} \\
greedy $2$-cover (1029E shape) & deepest-uncovered leaf $\to$ attach node~$1$ to its parent; on $T$ this adds $2$ edges ($(1,9)$, $(1,5)$) & \secref{c27-greedy} \\
Grundy on the stem & $\mathrm{mex}$/colon rule: $g=[7,3,2,0,1,0,2,0,1,0]$ for $v=1..10$, whole stalk $g(1)=7\neq0$ & \secref{c27-games} \\
\bottomrule
\end{tabular}
\end{center}

% ---------------- Taxonomy chart ----------------
\begin{figure}[H]
\centering
\begin{tikzpicture}[diagram, scale=0.95, transform shape,
    mnode/.style={circle, draw=#1, line width=0.9pt, fill=#1!8, minimum size=6.5mm,
                  font=\scriptsize\bfseries, text=#1!60!black}]
  \node[mnode=primaryblue] (a) at (0,0)     {$\mathit{dp}[v][j]$};
  \node[annot, font=\scriptsize, anchor=west] at (1.1,0)  {$=$ value over a budget $j=0..V$; merge children one at a time (14.1)};
  \node[mnode=successgreen] (b) at (0,-1)  {$\mathit{dp}[v][x]$};
  \node[annot, font=\scriptsize, anchor=west] at (1.1,-1) {$=$ profile indexed by a growing value/time/length (14.2)};
  \node[mnode=difforange]   (c) at (0,-2)  {$\pi(v)$};
  \node[annot, font=\scriptsize, anchor=west] at (1.1,-2) {$=$ greedy choice on a leaf-pruning order, proved by exchange (14.3)};
  \node[mnode=warningred]   (d) at (0,-3)  {$g(v)$};
  \node[annot, font=\scriptsize, anchor=west] at (1.1,-3) {$=$ Grundy $\mathrm{mex}$ of options, or a win/lose flag up the trie (14.4)};
  \node[mnode=softpurple]   (e) at (0,-4)  {$\mathit{ans}$};
  \node[annot, font=\scriptsize, anchor=west] at (1.1,-4) {$=$ read at the root, summed over splits, or $\bigoplus_v g(v)$ at the top};
\end{tikzpicture}
\caption*{The budgeted-DP taxonomy -- four state shapes, one per-child merge.
Rows 1--2 are the $\bigO{nV}$ family of \thmref{c27-amortize}; row~3 is a DP
that \emph{collapses} to a greedy because a tree has no cycles; row~4 is the
game value, a $\mathrm{mex}$ where the others are $\max$. The Ch.13
cookbook is the $V=\text{const}$, single-slot special case of row~1.}
\end{figure}

% ---------------- Who owns what ----------------
\subsection*{Who Owns What (so nothing is learned twice)}

\begin{center}
\small
\renewcommand{\arraystretch}{1.25}
\begin{tabular}{@{}p{0.30\linewidth} p{0.32\linewidth} p{0.32\linewidth}@{}}
\toprule
\textbf{Idea} & \textbf{Owner} & \textbf{This chapter's share} \\
\midrule
constant-tuple child-fold, in/out pair, harvest-at-peak, depth histogram & Ch.~13 (\chref{c10}) & the fold is now \emph{indexed by a budget}; Ch.~13's shapes are the $V=\text{const}$ special case \\
tree \emph{metrics} (depth, height, size, diameter, centroid, radius) & Ch.~12 (\chref{c09}) & treated as already-computed inputs; $\size(v)$ drives the loop cap \\
join-arms-at-the-peak for diameter & Ch.~12 (\thmref{c09-dpjoin}) & cited, not re-proved; the harvest is already Ch.~13's \\
sum-of-subtree-size / $\bigO{nV}$ child-merge bound & \textbf{Ch.~14 (this)} & 13.4 previewed it without proof; \thmref{c27-amortize} proves it here \\
small-to-large / DSU-on-tree merging & Ch.~17 (\chref{c28}) & \emph{not} here -- every merge is a plain per-child knapsack, never heavy-child reuse \\
rerooting as a theory / answer for every root & Ch.~15 (\chref{c20}) & \emph{not} here -- root is fixed throughout \\
per-pair distance \emph{queries} via LCA & Ch.~16 (\chref{c11}) & one-shot DP; queries live there \\
centroid decomposition for unbounded $k$ & Ch.~24 (\chref{c18}) & when the budget index is large, hand off \\
Cartesian / monotone-stack tree building & Ch.~(c26) & greedy here is \emph{on a given tree}, not constructing one \\
\bottomrule
\end{tabular}
\end{center}

% ---------------- Recognition guide ----------------
\begin{recognitionbox}[Pattern card -- ``the state has a numeric index and children compete for one budget'']
\reckey{the answer at a node is an \emph{array} $\mathit{dp}[v][0..V]$ over a
count/length/value budget; choosing how much budget each child gets is a knapsack;
or the statement is a two-player game on a tree/trie; or a covering/guarding
number where a leaf-sweep greedy is plausible.}
\recuse{
  \item ``at most $V$ special vertices in a subtree / pick a budget-capped subset''
        -- the tree knapsack of \secref{c27-knapsack}; bound the loops by
        \thmref{c27-amortize} (CF~855C, Summit CF~735E);
  \item ``the state is indexed by a value/time/length that ranges with the data''
        -- the profile DP of \secref{c27-growing} (CF~490F);
  \item ``cover/guard the tree, minimise guards/edges, no cycles to break the
        choice'' -- the exchange greedy of \secref{c27-greedy} (CF~1029E);
  \item ``two players, optimal play, last-move-or-unable-to-move loses'' on a
        tree or trie -- the win/lose or Grundy DP of \secref{c27-games} (CF~455B).
}
\recnotuse{
  \item the state is a \emph{constant} tuple (one number per node) -- that is
        Chapter~13's cookbook, read \chref{c10} first;
  \item you must reuse a \emph{large} per-node map and merge the smaller into
        the larger -- that is small-to-large / DSU-on-tree, Chapter~17 (\chref{c28});
  \item the same budgeted DP must be answered for \emph{every root} --
        rerooting, Chapter~15 (\chref{c20});
  \item the budget index $V$ is $\bigO{n}$ and you need $\bigO{n\log n}$ --
        the centroid-decomposition merge of Chapter~24 (\chref{c18}) takes over.
}
One-line test: \emph{does the node's answer depend on \underline{how much of a
fixed budget} you hand each child, forcing an array state merged child-by-child?
Then it is a budgeted tree DP, and this chapter is its home.}
\end{recognitionbox}

% ---------------- Constraint -> tool ----------------
\subsection*{Reading the Constraints (which budget the limits want)}

\begin{center}
\small
\renewcommand{\arraystretch}{1.25}
\begin{tabular}{@{}p{0.28\linewidth} p{0.40\linewidth} p{0.24\linewidth}@{}}
\toprule
\textbf{Constraint signature} & \textbf{What it rules in/out} & \textbf{Tool} \\
\midrule
$n\le10^{5}$, a count $x$ with $x\le10$ & array of size $x{+}1$ per node; per-child knapsack, amortised & \secref{c27-knapsack}, CF~855C \\
``at most $V$ of something, subtrees merge'' & bound inner loops by $\min(\text{so far},V)$ & \thmref{c27-amortize} \\
$n\le6000$, an increasing-sequence value along a path & $\bigO{n^2}$ profile / patience merge over paths & \secref{c27-growing}, CF~490F \\
cover within distance~$2$, minimise additions, $n\le2\cdot10^{5}$ & leaf-sweep greedy (a DP would TLE) & \secref{c27-greedy}, CF~1029E \\
two-player, total string length $\le10^{5}$, $k$ rounds & win/lose flag up a trie, or Grundy $\mathrm{mex}$ & \secref{c27-games}, CF~455B \\
$n\le100$, distance budget $k\le20$ on a tree knapsack & budgeted subtree knapsack under \thmref{c27-amortize} & Summit CF~735E \\
\bottomrule
\end{tabular}
\end{center}

% ---------------- Notation ----------------
\subsection*{Notation Used Only Here}

Root fixed at $1$ (or as stated); $c\in\mathit{ch}(v)$ ranges over the children
of $v$ in that rooted view; $\size(v)$ is the subtree size (Ch.~12) and it is
\emph{the running budget} -- after merging a child the "size so far" grows by
$\size(c)$. $\mathit{dp}[v][j]$ is a value over budget $j\in[0,V]$.
A \term{merge} of child $c$ into $v$ means
$\mathit{dp}[v][i+j]\!\leftarrow\!\mathit{dp}[v][i]+\mathit{dp}[c][j]$
for the admissible range of $i,j$ bounded by \thmref{c27-amortize}.
$\mathrm{mex}\,S$ is the least non-negative integer not in $S$; $g(v)$ is a
Grundy number; $\mathcal{G}=\bigoplus$ is XOR of component Grundy values.
A single leaf has $\mathit{dp}[\ell][0]=0$, $\mathit{dp}[\ell][j\ge1]$ from its
own pick, $g(\ell)=0$, and it is trivially "covered" in any distance-$\le V$ game.

Every section of this chapter is a chapter of the \emph{same} table on the
\emph{same} $T$: if a row here and a row there disagree, this opener is wrong
and the section is right -- fix it here, never the other way.
